\section{Continuous gradient lifts}
\label{sec:continuous-gradient}

There is a second way to pass from transport on the punctured zero set
to ambient right triviality.  It uses vector fields that extend
continuously to the origin section, together with a condition preventing
their trajectories from reaching that section in finite time.  The
gradient lift and its correction near the zero set belong to the
vector-field method of Kuo and Parusi\'nski; see
\cite[Section~2, especially equation~(2.5)]{Parusinski1999}.
We give the criterion with all its transport hypotheses explicit.

Let $G$ be holomorphic on a neighborhood of $\{0\}\times D$ in
$\mathbb C^n\times D$, where $D$ is a complex disk centered at zero,
and suppose $G(0,u)=0$.  Set
\[
 A=\{0\}\times D,\qquad X^*=\{G=0\}\setminus A.
\]
Real tangent vectors are written in complex coordinates as $(w,a)$,
with $w\in\mathbb C^n$ and $a\in\mathbb C$.  Norms are Euclidean.
All assertions below concern sufficiently small representatives.

\begin{proposition}[A continuous-lift criterion]
\label{prop:continuous-gradient-lift}
Assume that $d_xG\ne0$ off $A$ and that, at every point of $A$,
\begin{equation}
\label{eq:gradient-small-o}
 \frac{|G_u(x,u)|}{\|d_xG(x,u)\|}\longrightarrow0
 \qquad\text{as }(x,u)\longrightarrow A,\quad x\ne0.
\end{equation}
For each $a=1,i$, suppose that $X^*$ carries a smooth real vector field
$v_a=(w_a,a)$ with the following properties:
\begin{enumerate}
\item $dG(v_a)=0$ and $w_a(x,u)\to0$ on approach to every point of $A$;
\item no trajectory of $v_a$ in $X^*$ has a finite-time endpoint in
      $A$ lying inside the representative, in either time direction.
\end{enumerate}
Then, after shrinking, there is an ambient homeomorphism germ
\[
 H(x,u)=(H_u(x),u),\qquad
 G(H_u(x),u)=G(x,0),\qquad H_u(0)=0,
\]
with $H_0=\mathrm{id}$ and with $H$ and $H^{-1}$ jointly continuous.
\end{proposition}

\begin{proof}
On the ambient complement of $A$, put
\[
 k_a=-aG_u\,
       \frac{(\overline{G_{x_1}},\ldots,\overline{G_{x_n}})}
            {\|d_xG\|^2},\qquad K_a=(k_a,a).
\]
These smooth real fields satisfy $dG(K_a)=0$.
Equation~\eqref{eq:gradient-small-o} says precisely that $k_a\to0$ at
$A$.  On $X^*$, the difference $s_a=w_a-k_a$ is a section of the
complex vector bundle $K=\ker d_xG$, and it too tends to zero at $A$.

We extend this correction while preserving its vanishing at the
section.  The smooth submanifold $X^*$ is closed in the ambient
complement of $A$.  Choose a smooth tubular retraction $r:U\to X^*$,
with a variable tube radius small enough that
\[
 \|z-r(z)\|<\tfrac12\operatorname{dist}(r(z),A)
 \qquad(z\in U).
\]
Choose a smooth function $\chi$ on the axis complement, equal to one
on $X^*$, taking values in $[0,1]$, and with support in $U$.  If
$\pi_{K_z}$ denotes Hermitian orthogonal projection onto $K_z$, define
\[
 \widetilde s_a(z)=\chi(z)\pi_{K_z}(s_a(r(z)))\quad(z\in U),
 \qquad \widetilde s_a=0\quad\text{outside }U.
\]
This is a smooth section of $K$ agreeing with $s_a$ on $X^*$.
If $z$ approaches a point of $A$ within $U$, the displayed tube bound
implies that $r(z)$ approaches the same point.  Since the projection
has norm at most one, $\widetilde s_a(z)\to0$ as well.  Thus
\[
 W_a=(k_a+\widetilde s_a,a)
\]
extends continuously across $A$ by $W_a=(0,a)$.  It is smooth off
$A$, agrees with $v_a$ on $X^*$, and annihilates $dG$ everywhere.
At points of $A$ the last assertion follows from $G(0,u)=0$.

Peano's existence theorem gives local solutions of each real ODE
$\dot z=W_a(z)$.  They are unique off $A$ because the field is
smooth there.  A solution meeting $A$ has $G=0$ throughout, by the
chain rule.  If it left $A$, a connected interval on which it lies
off $A$ would give a trajectory of $v_a$ with a finite-time endpoint
on $A$, contrary to the hypothesis.  The same argument in reverse
time excludes arrival from outside $A$.  On $A$ the equation is
$\dot x=0$, $\dot u=a$, so uniqueness holds there too.

For clarity, continuity of these local flows also holds at $A$.
On a fixed compact neighborhood the continuous field is bounded.
Initial points in a smaller neighborhood therefore have solutions on
a common short interval, all contained in that compact neighborhood.
For a convergent sequence of initial points, the solutions are
uniformly bounded and equicontinuous.  Every uniformly convergent
subsequence satisfies the integral equation for the limiting initial
point.  Uniqueness identifies its limit, and hence the whole sequence
converges.  The same argument applies to negative times.  Uniqueness
and reverse flow give continuous local inverse maps.  Consequently
these are local flows by homeomorphisms, jointly continuous in the
initial point and time.

Write $u=s+it$.  Starting at $(x,0)$, first follow $W_1$ for time
$s$ and then $W_i$ for time $t$.  Uniform short-time existence permits
both steps for all $x$ and $u$ in common smaller neighborhoods.  The
endpoint has parameter $u$ and defines $H_u(x)$.  Each step preserves
$G$ and the section.  Reverse the two steps, in the opposite order,
to obtain the joint inverse.  The reverse operations are defined on
an open neighborhood of the section.  Within that neighborhood the
image of our open source is its inverse image under the reverse map,
so this image is open.  At $u=0$ both times are zero.  This
proves all the assertions.
\end{proof}

\begin{remark}[Radial transport]
\label{rem:continuous-radial}
If the no-finite-time-endpoint condition holds for every constant real
combination of $v_1$ and $v_i$, one may instead integrate
$\operatorname{Re}(u)W_1+\operatorname{Im}(u)W_i$ for time one.
The preceding uniqueness argument applies to this combined field.
Its speed on a compact neighborhood is bounded by a constant times
$|u|$, so the solutions exist up to time one after shrinking $u$ and
the initial neighborhood.  Applying the same compactness argument to
the integral equations, now also allowing $u$ to vary, gives joint
continuity.  Reversing the trajectory gives the inverse.  Thus this
stronger trajectory hypothesis provides transport along $[0,u]$.
\end{remark}

The small-o condition and the trajectory condition have different
roles.  The L\^e--Saito--Teissier criterion \cite{LeSaito1973} supplies
\eqref{eq:gradient-small-o} for the $\mu$-constant isolated families
considered here; see \cite[Chapter~II, Remark~3.10]{Teissier1973} and
\cite[proof of Corollary~2.1, equation~(2.6)]{Parusinski1999}.
The latter writes the estimate with the full gradient; it is
equivalent to the spatial-gradient version since
$\|dG\|^2=\|d_xG\|^2+|G_u|^2$.
The following two constructions supply the additional trajectory
condition.

\begin{corollary}[An explicit Whitney lift]
\label{cor:gradient-whitney}
Suppose that $G$ satisfies the analytic and gradient hypotheses of
Proposition~\ref{prop:continuous-gradient-lift} and that $(X^*,A)$
satisfies the Whitney conditions.  Then $G$ has the trivialization
of that proposition.  The fields on $X^*$ can be chosen to preserve
$\|x\|^2$, and every constant real combination has the same property.
\end{corollary}

\begin{proof}
Use the Hermitian inner product linear in its first entry, and set
$q(x,u)=\pi_{K_{(x,u)}}x$.  We first show that
\begin{equation}
\label{eq:whitney-projection}
 \frac{\|q(x,u)\|}{\|x\|}\longrightarrow1
 \qquad\text{on }X^*\text{ approaching }A.
\end{equation}
Along any sequence, pass to a subsequence on which $x/\|x\|$ and
$K_{(x,u)}$ converge, to $v$ and $Q$ respectively.
Equation~\eqref{eq:gradient-small-o} implies that the tangent spaces
of the total zero set converge to $Q\times\mathbb C$.
Whitney condition~(b), applied to the secants between $(x,u)$ and
$(0,u)$, puts $(v,0)$ in this limiting space.  Thus $v\in Q$, which
proves \eqref{eq:whitney-projection}.

In particular $q\ne0$ sufficiently near the section.  With $k_a$
as in the preceding proof, define
\[
 w_a=k_a-
      \frac{\operatorname{Re}\langle k_a,x\rangle}{\|q\|^2}\,q,
 \qquad v_a=(w_a,a),\qquad a=1,i.
\]
The correction belongs to $K$, so $dG(v_a)=0$.
Moreover $\operatorname{Re}\langle q,x\rangle=\|q\|^2$, giving
$d\|x\|^2(v_a)=0$.  Its norm is bounded by
$\|k_a\|\|x\|/\|q\|$; hence $w_a\to0$ by
\eqref{eq:gradient-small-o} and \eqref{eq:whitney-projection}.
Every constant real combination preserves the positive radius of a
trajectory in $X^*$, so no such trajectory can reach $A$ in finite
time.  Proposition~\ref{prop:continuous-gradient-lift} and
Remark~\ref{rem:continuous-radial} apply.
\end{proof}

\begin{corollary}[Continuous lifts from a resolution]
\label{cor:gradient-resolution}
Suppose that $G$ satisfies the analytic and gradient hypotheses of
Proposition~\ref{prop:continuous-gradient-lift}.  Suppose that there
is a proper holomorphic map $b:M\to\{G=0\}$ over $D$, an
isomorphism away from $A$, such that the composite $q:M\to D$ is
smooth and $E=b^{-1}(A)$ is a relative
normal-crossing divisor.  Then $G$ has the trivialization of that
proposition.  The fields on $X^*$ satisfy the trajectory condition
for every constant real combination.
\end{corollary}

\begin{proof}
Relative normal-crossing coordinates give smooth lifts of the real
parameter directions $1,i$ tangent to every local branch of $E$.
Submersion coordinates supply lifts away from $E$.  A partition of
unity patches these to smooth fields $\widetilde v_1,\widetilde v_i$
on $M$ with the same projections and tangencies.  Through the
isomorphism off $E$, push them to fields $v_a=(w_a,a)$ on $X^*$.

These spatial parts tend to zero at $A$.  Indeed, properness supplies
convergent subsequences of the unique preimages of any sequence in
$X^*$ tending to a point of $A$.  Every limit lies in $E$.  The
spatial component of $b$ vanishes on $E$, so its differential on a
field tangent to a branch of $E$ is zero there.  Continuity of
$db(\widetilde v_a)$ now proves $w_a\to0$.

To verify the trajectory condition, suppose that an integral arc of
a constant real combination of the $v_a$ had a finite-time endpoint
in $A$ inside the representative.  Its lift to $M\setminus E$
solves the corresponding smooth combined field.  A terminal segment
of the downstairs arc has compact closure, whose inverse image is
compact by properness.  Smooth ODE continuation therefore extends
the lifted arc to a point of $E$.  Tangency preserves $E$, and
uniqueness in both time directions prevents an arc outside $E$ from
reaching it.  This contradiction proves the required condition.
The continuous-lift criterion applies.
\end{proof}

For the families in the main theorem, these criteria give an alternative
construction of the ambient maps.  In the double-point case of
Section~\ref{sec:low-cases}, apply the continuous Whitney criterion to
the plane family $g_t$ after splitting off a square, and then extend
the maps by the identity in the remaining coordinate and conjugate back
as in that section.
The log canonical case of multiplicity three uses the Whitney criterion
directly.  In the non-log-canonical case, use the resolved family of
Proposition~\ref{prop:resolved-family} after its finite disk cover.
The L\^e--Saito--Teissier estimate applies to each of these
$\mu$-constant families.  The closed-ball construction in
Section~\ref{sec:ambient-lifting} also supplies exact common tube
domains, which are used in Lemma~\ref{lem:remove-base-change} to
descend from the finite cover to the original disk.
